
\chapter{Visão Computacional}
\label{chap:visao-computacional}

% Apresentar o papel da visão computacional no DoseCerta
% e a função que ela desempenha no sistema.

\section{Objetivo e Requisitos}
\label{sec:objetivo-visao-computacional}

% Descrever o problema que a visão computacional busca resolver.
% Relacionar a função aos requisitos do sistema e ao público-alvo.

\section{Dispositivo de Captura de Imagens}
\label{sec:captura-imagens}

% Identificar a câmera ou dispositivo utilizado.
% Informar características relevantes, como resolução e condições
% de captura, conforme documentação ou configuração real.

\section{Conjunto de Dados}
\label{sec:conjunto-dados-visao}

% Descrever as imagens utilizadas, sua origem e organização.
% Se houve coleta própria, explicar como foi realizada.
% Informar critérios de seleção, divisão e anotação dos dados,
% caso aplicáveis. Não declarar treinamento se não ocorreu.

\section{Pré-processamento das Imagens}
\label{sec:preprocessamento-imagens}

% Explicar as operações aplicadas às imagens antes da análise,
% como redimensionamento, recorte ou normalização,
% somente se utilizadas na implementação.

\section{Método de Análise}
\label{sec:metodo-analise-imagens}

% Apresentar o algoritmo, modelo ou técnica adotada.
% Explicar sua finalidade e o princípio de funcionamento
% com apoio em referências, quando necessário.

\section{Implementação}
\label{sec:implementacao-visao-computacional}

% Descrever as ferramentas, bibliotecas e ambiente utilizados.
% Explicar o fluxo de processamento implementado e suas entradas
% e saídas. Inserir pseudocódigo ou diagrama, se pertinente.

\section{Integração com o DoseCerta}
\label{sec:integracao-visao-computacional}

% Explicar como o resultado da análise é utilizado pelo sistema.
% Descrever a comunicação com o backend, aplicativo ou sistema
% embarcado, conforme a arquitetura implementada.

\section{Limitações e Condições de Uso}
\label{sec:limitacoes-visao-computacional}

% Registrar limitações conhecidas, como condições de iluminação,
% posicionamento, qualidade da imagem ou casos não reconhecidos.
% Diferenciar limitações observadas de riscos ainda não avaliados.

\section{Síntese da Visão Computacional}
\label{sec:sintese-visao-computacional}

% Resumir a abordagem adotada, sua contribuição ao DoseCerta
% e o estágio atual de desenvolvimento.
